\documentclass[11pt, a4paper]{article}

% --- PACKAGES ---
\usepackage{graphicx}
\usepackage{amsmath}
\usepackage{amssymb}
\usepackage{geometry}
\usepackage{booktabs}
\usepackage{hyperref}
\usepackage{listings}
\usepackage{xcolor}
\usepackage{float}

% --- CONFIG ---
\geometry{margin=1in}
\definecolor{codegray}{rgb}{0.95,0.95,0.95}
\lstset{
    backgroundcolor=\color{codegray},
    basicstyle=\ttfamily\small,
    breaklines=true,
    frame=single,
    columns=fullflexible
}

% --- TITLE ---
\title{\textbf{Project RIFT: Memory as Irreversible Topological Deformation}\\
\large Failure Modes of Learning in a Non-Equilibrium Antagonistic Substrate}

\author{\textbf{Akshay John} \\ Independent Researcher \\ \texttt{Experiment Series: RIFT-003}}
\date{January 2026}

\begin{document}
\maketitle

% --- ABSTRACT ---
\begin{abstract}
Most artificial intelligence systems rely on minimizing an externally defined loss function. We explore an alternative paradigm in which adaptation emerges from physical necessity rather than optimization. We present \textbf{RIFT (Resilient Internal Fracture Topology)}, a non-equilibrium substrate governed by antagonistic homeostatic regulation. The system does not solve tasks or optimize objectives; instead, it must continuously reorganize its internal structure to avoid thermodynamic collapse or runaway instability.

Through a sequence of controlled failures and refinements, we demonstrate that memory can emerge as \textbf{irreversible topological deformation} (“scarring”) without backpropagation, targets, or reward signals. We further show that memory capacity is severely constrained by global homeostatic dynamics, leading to insulation and catastrophic interference. These results suggest that learning in self-regulating substrates requires separation of sensing, survival, and consolidation mechanisms operating on distinct physical timescales.
\end{abstract}

% --- INTRODUCTION ---
\section{Introduction}

Contemporary AI architectures are fundamentally episodic: computation occurs only when externally invoked, and internal state does not persist unless explicitly stored. In contrast, biological systems exist in a continuous struggle against entropy, where structure is maintained through constant energy dissipation.

Project RIFT investigates whether adaptive behavior can emerge from a system that has no task, no reward, and no objective—only the requirement to persist under antagonistic physical constraints. Rather than asking whether the system can perform intelligence, we ask a more primitive question: \textit{what physical conditions are required for memory to exist at all?}

This work does not claim to implement intelligence. Instead, it studies the failure modes that arise when survival, learning, and structural evolution are forced to coexist within a single substrate.

% --- ARCHITECTURE ---
\section{System Architecture: Antagonistic Homeostasis}

The RIFT system consists of three interacting components:

\begin{enumerate}
    \item \textbf{Tectonic Substrate:} A $48 \times 48$ spatial grid with 32 channels. Local interactions diffuse energy and propagate stress.
    \item \textbf{Fracture Law:} If local stress exceeds elastic resistance, the substrate fractures. Fracture dissipates energy and permanently increases local resistance.
    \item \textbf{Antagonist EGO:} A global regulator enforcing non-equilibrium:
    \begin{itemize}
        \item If fracture rate is too low (stagnation), random perturbations are injected.
        \item If fracture rate is too high (instability), energy is dampened.
    \end{itemize}
\end{enumerate}

No gradients, targets, rewards, or optimization signals are used.

% --- DEVELOPMENT PHASES ---
\section{Developmental Phases and Early Failures}

\subsection{Phase I: Unbounded Dynamics}
Without dissipation limits, fracture amplified stress, leading to numerical divergence:
\begin{lstlisting}
Step 030 | Fracture=2.8e+23 -> EXPLOSION
\end{lstlisting}

\subsection{Phase II: Over-Damping}
Heavy damping eliminated instability but led to total stagnation:
\begin{lstlisting}
Step 450 | Fracture=0.0000 -> COMA
\end{lstlisting}

\subsection{Phase III: Antagonistic Regulation}
Introducing the Antagonist EGO forced sustained non-equilibrium dynamics:
\begin{lstlisting}
Step 025 | Fracture=0.0019 | PRESSURIZE
Step 075 | Fracture=0.6858 | CONSTRAIN
\end{lstlisting}

This produced continuous oscillation without collapse.

% --- MEMORY EMERGENCE ---
\section{Experiment I: Episodic Trauma and Physical Memory}

\subsection{Protocol}
A high-magnitude impulse (“Hammer Strike”) was applied at fixed spatial locations every 50 timesteps. Elasticity decay was minimized to preserve long-term effects.

\subsection{Result}
Peak fracture at impact decreased monotonically across repetitions:
\begin{lstlisting}
0.0644 → 0.0271 → 0.0092 → 0.0036
\end{lstlisting}

\subsection{Interpretation}
The substrate physically reorganized to resist a specific stress pattern. This constitutes memory as irreversible deformation rather than representational encoding.

% --- CAPACITY TESTS ---
\section{Experiment II: Spatial Memory Capacity (MNEMOS-II / III)}

\subsection{Protocol}
Nine spatial locations were trained sequentially using identical trauma. After training, each location was struck once to assess recall.

\subsection{Result}
Retention degraded rapidly:
\begin{itemize}
    \item MNEMOS-II: 1/9 sites retained
    \item MNEMOS-III (tuned): 4/9 sites retained
\end{itemize}

\subsection{Failure Mode: Global Callus}
Background perturbations from the Antagonist EGO caused diffuse micro-fracturing, incrementally hardening the entire substrate. Later sites failed to fracture at all, preventing learning.

% --- DUAL TIMESCALE ---
\section{Experiment III: Dual-Timescale Elasticity (MNEMOS-IV)}

\subsection{Motivation}
To decouple survival from learning, elasticity was split into:
\begin{itemize}
    \item Fast elasticity (rapid adaptation, rapid decay)
    \item Slow elasticity (consolidated memory)
\end{itemize}

\subsection{Result}
Despite protecting early memories, later sites remained pre-hardened:
\begin{itemize}
    \item Retained sites: 2/9
\end{itemize}

\section{Experiment III: Spatially Gated Dual-Timescale Memory}

\subsection{Motivation}
Previous experiments (MNEMOS-IV) revealed that dual timescales alone are insufficient; homeostatic regulation causes global pre-hardening ("Insulation"), blocking future memory formation. To address this, we introduced a **Spatial Gating Mechanism** (Novelty Filter).

\subsection{Method}
We implemented a local sharpness filter based on the Laplacian of the stress field:
\[ S(x,y) = | \sigma(x,y) - \mu_{local} | \]
Learning in the slow-decay memory layer ($\epsilon_{slow}$) is gated by $S(x,y)$. This allows the substrate to differentiate between diffuse homeostatic noise (low sharpness) and structural trauma (high sharpness).

\subsection{Results}
The Spatially Gated substrate successfully retained \textbf{7 out of 9} memory sites ($78\%$ Capacity).
\begin{itemize}
    \item \textbf{Retention:} Sites $X_1, X_3, X_4, X_5, X_7, X_8, X_9$ showed clear episodic learning ($Fracture_{recall} \ll Fracture_{initial}$).
    \item \textbf{Interference:} Sites $X_2$ and $X_6$ were pre-hardened due to stress diffusion from neighboring sites, confirming the physical (non-symbolic) nature of the substrate.
\end{itemize}

\subsection{Conclusion}
We have demonstrated that a non-optimizing, antagonistically regulated substrate can support high-capacity, spatially distributed memory. By decoupling \textit{survival dynamics} (Fast Skin) from \textit{memory consolidation} (Slow Bone) via spatial gating, the system resolves the Stability-Plasticity dilemma without gradient descent.

\subsection{Key Insight}
Separating learning rates is insufficient. Because fracture detection depends on total resistance, fast-layer hardening globally suppresses future damage detection, blocking consolidation.

% --- FAILURE ANALYSIS ---
\section{Failure Modes and Constraints}

\subsection{Global Insulation}
Homeostatic mechanisms raise resistance everywhere, preventing new experiences from registering as damage.

\subsection{Stability–Plasticity Conflict}
Survival dynamics overwrite learning capacity even in non-optimizing systems.

\subsection{Structural Limitation}
A single damage variable cannot simultaneously support sensing, survival, and memory.

% --- SUMMARY TABLE ---
\begin{table}[H]
\centering
\begin{tabular}{lccc}
\toprule
Experiment & Objective & Outcome & Verdict \\
\midrule
RIFT v3 & Sustained Non-Equilibrium & Stable Oscillation & Pass \\
Episodic Trauma (v4-B) & Physical Memory & Fracture ↓ 18× & Pass \\
Mutation (v4-C) & Structural Evolution & Memory Destroyed & Expected \\
MNEMOS-II & Capacity Test & 1 / 9 Retained & Fail \\
MNEMOS-III & Tuned Capacity & 4 / 9 Retained & Partial \\
MNEMOS-IV & Dual Timescale & 2 / 9 Retained & Fail \\
\bottomrule
\end{tabular}
\caption{Summary of Experimental Regimes}
\end{table}

% --- CONCLUSION ---
\section{Conclusion and Future Work}

Project RIFT demonstrates that memory can emerge from irreversible physical deformation without optimization, targets, or symbolic representation. However, it also reveals a fundamental limitation: systems that must actively survive tend to insulate themselves against novelty, suppressing future learning.

Future work (MNEMOS-V) will investigate separating damage \textit{sensing} from damage \textit{response}, introducing predictive stress detection prior to elastic hardening. This separation may be a necessary precondition for scalable learning in self-regulating substrates.

\end{document}
